\documentclass{scrartcl}
\usepackage[utf8]{inputenc}
\usepackage{lmodern} \usepackage{amsmath}
\usepackage{amssymb}
\usepackage{amsthm}
\usepackage{thmtools}
\usepackage{mathtools}
\usepackage{hyperref} \usepackage{enumitem} \usepackage{graphicx} \usepackage{dsfont} \usepackage{tikz} \usepackage[framemethod=TikZ]{mdframed}
\usepackage{ifthen}
\usepackage{supertabular}  \newcommand*{\todo}[1]{[\textbf{\color{red} Todo:} \textit{#1}]}
\newcommand*{\note}[1]{[\textbf{\color{blue} Note:} \textit{#1}]}
\newenvironment{mynote}{\begin{quote}\color{gray}\textbf{\underline{Note:} }}{\end{quote}}
\renewcommand{\subset}{\subseteq}
\newcommand{\lcb}{\left\lbrace} \newcommand{\rcb}{\right\rbrace} \newcommand{\cb}[1]{\lcb #1 \rcb} \newcommand{\cbOf}[1]{\mathopen{}\lcb #1 \rcb\mathclose{}} \newcommand{\lab}{\left[} \newcommand{\rab}{\right]} \newcommand{\ab}[1]{\lab #1 \rab} \newcommand{\abOf}[1]{\!\ab{#1}} \newcommand{\lb}{\left(} \newcommand{\rb}{\right)} \newcommand{\br}[1]{\lb #1 \rb} \newcommand{\brOf}[1]{\!\br{#1}} \newcommand{\abs}[1]{\left| #1 \right|} \newcommand*{\E}{\mathbb{E}} \newcommand*{\ent}{\mathsf{Ent}} \newcommand*{\V}{\mathbb{V}} \newcommand*{\cov}{\mathbb{COV}} \let\Pr\relax \newcommand*{\Pr}{\mathbb{P}} \newcommand*{\powset}{\mathfrak{P}} \newcommand{\sizedMid}[2]{#1 \, \kern-\nulldelimiterspace\mathopen{}\left| \vphantom{#1}\,#2\right.\mathclose{}\kern-\nulldelimiterspace}
\newcommand{\EOf}[1]{\E\abOf{#1}}
\newcommand{\Eof}[1]{\E[#1]}
\newcommand{\entOf}[1]{\ent\brOf{#1}}
\newcommand{\entof}[1]{\ent(#1)}
\newcommand{\Vof}[1]{\V[#1]}
\newcommand{\VOf}[1]{\V\abOf{#1}}
\newcommand{\covof}[1]{\cov[#1]}
\newcommand{\ECondOf}[2]{\EOf{\sizedMid{#1}{#2}}}
\newcommand{\ECondof}[2]{\E[#1\mid #2]}
\newcommand{\PrOf}[1]{\Pr\mathopen{}\lb #1 \rb\mathclose{}}
\newcommand{\Prof}[1]{\Pr(#1)}
\newcommand{\PrCondOf}[2]{\PrOf{\sizedMid{#1}{#2}}}
\newcommand{\PrCondof}[2]{\Pr(#1\mid #2)}
\newcommand{\powsetOf}[1]{\powset\mathopen{}\lb #1 \rb\mathclose{}}
\newcommand{\eventAnd}{,\,} \newcommand{\setByEle}[2]{\cb{\sizedMid{#1}{#2}}}
\newcommand{\setByEleInText}[2]{\{#1 \mid #2\}}
\DeclareMathOperator{\diam}{\mathsf{diam}}
\DeclareMathOperator{\card}{\mathsf{card}}
\DeclareMathOperator{\ball}{\mathrm{B}}
\newcommand{\id}{\mathrm{id}}
\newcommand{\restrict}[1]{_{\mkern 1mu \vrule height 2ex\mkern2mu {#1}}}
\providecommand\given{} \newcommand\SetSymbol[1][]{
	\nonscript\,#1\vert \allowbreak \nonscript\,\mathopen{}}
\DeclarePairedDelimiterX\Set[1]{\lbrace}{\rbrace}{ \renewcommand\given{\SetSymbol[\delimsize]} #1 }
\newcommand{\Ex}{\E\expectarg}
\DeclarePairedDelimiterX{\expectarg}[1]{[}{]}{\ifnum\currentgrouptype=16 \else\begingroup\fi
	\activatebar#1
	\ifnum\currentgrouptype=16 \else\endgroup\fi
}
\newcommand{\innermid}{\nonscript\;\delimsize\vert\nonscript\;}
\newcommand{\activatebar}{\begingroup\lccode`\~=`\|
	\lowercase{\endgroup\let~}\innermid 
	\mathcode`|=\string"8000
}
\newcommand{\cardOf}[1]{\left\vert{#1}\right\vert}
\newcommand{\cardof}[1]{\vert{#1}\vert}
\newcommand*{\im}{\mathrm{im}}
\newcommand*{\mc}[1]{\mathcal{#1}}
\newcommand*{\mb}[1]{\mathbb{#1}}
\newcommand*{\mr}[1]{\mathrm{#1}}
\newcommand*{\ms}[1]{\mathsf{#1}}
\newcommand*{\mo}[1]{\mathbf{#1}}
\newcommand*{\mf}[1]{\mathfrak{#1}}
\newcommand{\N}{\mathbb{N}}
\newcommand{\Nn}{\mathbb{N}_0}
\newcommand{\R}{\mathbb{R}}
\newcommand{\Rp}{\R_+} \newcommand{\Ra}{\bar\R} \newcommand{\Q}{\mathbb{Q}}
\newcommand{\Z}{\mathbb{Z}}
\DeclareMathOperator{\borel}{\mathcal{B}}
\newcommand{\borelof}[1]{\borel(#1)}
\newcommand{\C}[2]{\mathcal{C}^{#1}\br{#2}}
\newcommand{\derive}[1]{#1^\prime}
\newcommand{\transpose}{\!^\top\!}
\newcommand{\tr}{\transpose}
\newcommand{\pr}{^\prime}
\newcommand{\prr}{^{\prime\prime}}
\newcommand{\prrr}{^{\prime\prime\prime}}
\newcommand{\normIntervall}{\lab 0, 1 \rab} \def\integral from #1to #2of #3by #4;{\int_{#1}^{#2} \! #3 \mathrm{d}#4} \def\integralMeasure in #1of #2by #3of #4;{\int_{#1} \! #2{#4} #3{\mathrm{d}#4}} \def\mapping #1from #2to #3;{#1 \colon #2 \rightarrow #3}
\def\mappingDef #1from #2to #3maps #4to #5;{#1 \colon #2 \rightarrow #3,\ #4 \mapsto #5}
\def\seq #1by #2;{\br{#1}_{#2\in\N}}
\def\seqInText #1by #2;{(#1)_{#2\in\N}}
\newcommand{\innerProduct}[2]{\left\langle#1\,,\, #2\right\rangle}
\newcommand{\ip}[2]{\innerProduct{#1}{#2}}
\newcommand{\lebesgue}{\mathcal{L}}
\newcommand{\lebesguePow}[1]{\lebesgue^{#1}}
\newcommand{\lebesgueOf}[1]{\lebesgue\brOf{#1}}
\newcommand{\invert}[1]{#1^{-1}}
\newcommand{\sgn}{\mathsf{sgn}} \newcommand{\supp}{\overline{\mathsf{supp}}} \newcommand{\Cut}{\mathsf{Cut}} \newcommand{\median}{\mathsf{median}} \newcommand{\Exp}{\mathsf{Exp}}
\newcommand{\Log}{\mathsf{Log}}
\newcommand{\dl}{\mathrm{d}}
\newcommand\independent{\protect\mathpalette{\protect\independenT}{\perp}}
\def\independenT#1#2{\mathrel{\rlap{$#1#2$}\mkern2mu{#1#2}}}
\def\converges for #1to #2;{\xrightarrow{#1} #2}
\def\convergesAlmostSurely for #1to #2;{\xrightarrow{#1}_{\mathsf{fs}} #2}
\def\convergesInProbability for #1to #2;{\xrightarrow{#1}_{\mathsf{p}} #2}
\def\convergesInL #1for #2to #3;{\xrightarrow{#2}_{\lebesguePow{#1}} #3}
\newcommand{\as}{\mathsf{a.s.}}
\newcommand{\ind}{\mathds{1}}\newcommand{\indOf}[1]{\ind_{\!#1}}\newcommand{\indOfOf}[2]{\ind_{\!#1}\!\brOf{#2}}\newcommand{\indOfEvent}[1]{\indOf{\cbOf{#1}}}\newcommand{\indOfEventOf}[2]{\indOf{\cbOf{#1}}\!\brOf{#2}}\newcommand{\compl}{^\mathsf{c}}\DeclareMathOperator{\landauO}{\mathcal{O}}\newcommand{\landauOOf}[1]{\landauO\brOf{#1}}\DeclareMathOperator{\landauOmega}{\Omega}\newcommand{\landauOmegaOf}[1]{\landauOmega\brOf{#1}}\DeclareMathOperator{\landauTheta}{\Theta}\newcommand{\landauThetaOf}[1]{\landauTheta\brOf{#1}}\newcommand{\norm}{\left\Vert \cdot \right\Vert}
\newcommand{\normof}[1]{\Vert #1 \Vert}
\newcommand{\normOf}[1]{\left\Vert #1 \right\Vert}
\newcommand{\unif}[1]{\mathsf{Unif}\brOf{#1}}
\newcommand{\equationFullstop}{\, .}
\newcommand{\eqfs}{\equationFullstop}
\newcommand{\equationComma}{\, ,}
\newcommand{\eqcm}{\equationComma}
\newcommand{\eqand}{\qquad\text{and}\qquad}
\newcommand{\sigmaAlg}{\mathrm{\sigma}}
\newcommand{\sigmaAlgof}[1]{\sigmaAlg(#1)}
\newcommand{\sigmaAlgOf}[1]{\sigmaAlg\brOf{#1}}
\newcommand{\euler}{\mathrm{e}}
\newcommand{\T}{^{T}}
\DeclareMathOperator*{\argmin}{arg\,min}
\DeclareMathOperator*{\argmax}{arg\,max}
\DeclareMathOperator*{\outerlim}{lim\, \overline{sup}}
\DeclareMathOperator*{\innerlim}{innerlim}
\newcommand*{\evmin}{\lambda^{\mathsf{min}}}
\newcommand*{\evmax}{\lambda^{\mathsf{max}}}
\newcommand{\entInt}[1]{\mathrm J[#1]}
\def\MS{\mathcal{Q}}\def\epi{\mathsf{epi}}\def\aec{\mathsf{aec}}\def\alsec{\mathsf{alsec}}\def\aecp{\mathsf{aec}_P}\def\Hausdorff{\mathsf{H}}\def\Haus{\mathsf{H}}\def\ubs{\mathsf{ubs}}\newcommand{\epiconv}[1]{\xrightarrow{#1\to\infty}_\epi}\newcommand{\hausconv}[1]{\xrightarrow{#1\to\infty}_\Hausdorff}\newcommand{\ubsconv}[1]{\xrightarrow{#1\to\infty}_\ubs}\newcommand{\ol}[2]{\overline{#1,\!#2}}
\newcommand{\sol}[2]{\overline{#1,#2}}
 \theoremstyle{plain}\newtheorem{theorem}{Theorem}[section]
\newtheorem{lemma}[theorem]{Lemma}
\newtheorem{corollary}[theorem]{Corollary}
\newtheorem{proposition}[theorem]{Proposition}
\theoremstyle{definition}
\newtheorem{definition}[theorem]{Definition}
\newtheorem{example}[theorem]{Example}
\theoremstyle{remark}
\newtheorem{remark}[theorem]{Remark}
\newtheorem{notation}[theorem]{Notation}
\newtheorem{assumptions}[theorem]{Assumptions}
\begin{document}
\title{Strong Laws of Large Numbers for Generalizations of Fréchet Mean Sets}
\subtitle{\url{https://github.com/chroetz/PaperStrong22}}
\author{Christof Schötz\\\href{mailto:math@christof-schoetz.de}{math@christof-schoetz.de}}
\date{}
\maketitle
\begin{abstract} 
A Fréchet mean of a random variable $Y$ with values in a metric space $(\mathcal Q, d)$ is an element of the metric space that minimizes $q \mapsto \mathbb E[d(Y,q)^2]$. This minimizer may be non-unique. 
We study strong laws of large numbers for sets of generalized Fréchet means.
Following generalizations are considered: 
the minimizers of $\mathbb E[d(Y, q)^\alpha]$ for $\alpha > 0$,
the minimizers of $\mathbb E[H(d(Y, q))]$ for integrals $H$ of non-decreasing functions, and
the minimizers of $\mathbb E[\mathfrak c(Y, q)]$ for a quite unrestricted class of cost functions $\mf c$.
We show convergence of empirical versions of these sets in outer limit and in one-sided Hausdorff distance. The derived results require only minimal assumptions.
 \end{abstract}
\section{Fréchet Mean Sets}
For a random variable $Y$ with values in $\R^s$ and $\Eof{\normof{Y}^2}<\infty$, it holds
\begin{equation*}
	\Ex{Y} = \argmin_{q\in\R^s} \Ex{d(Y,q)^2}
	\eqcm
\end{equation*}
where $d(x,y) = \normof{x-y}$ is the Euclidean distance. We can also write 
\begin{equation*}
	\Ex{Y} = \argmin_{q\in\R^s} \Ex{d(Y,q)^2-d(Y,0)^2}
	\eqcm
\end{equation*}
as we just add a constant term. In the latter equation, we only require $Y$ to be once integrable, $\Eof{\normof{Y}}<\infty$, instead of twice as $|\normof{y-q}^2-\normof{y}^2| \leq 2\normof{y}\normof{q}+\normof{q}^2$.

The concept of Fréchet mean, proposed in \cite{frechet48}, builds upon this minimizing property of the Euclidean mean to generalize the expected value to random variables with values in a metric space. Let $(\mc Q, d)$ be a metric space. As a shorthand we may write $\ol qp$ instead of $d(q,p)$. Let $Y$ be a random variable with values in $\mc Q$. Fix an arbitrary element $o\in \mc Q$. The \emph{Fréchet mean set} of $Y$ is $M = \argmin_{q\in\mc Q} \Ex{\ol Yq^2- \ol Yo^2}$ assuming the expectations exist. 
This definition does not depend on $o$. The reason for subtracting $\ol Yo^2$ is the same as in the Euclidean case: We need to make less moment assumptions to obtain a meaningful value: The triangle inequality implies
\begin{equation*}
\abs{\ol yq^2 - \ol yo^2} \leq \ol oq \br{\ol oq + 2 \ol yo}
\end{equation*}
for all $y,q,o\in\mc Q$. Thus, if $\Ex{\ol Yo} < \infty$, then $\Ex{\ol Yq^2 - \ol Yo^2} < \infty$ for all $q\in\mc Q$.

In Euclidean spaces and other Hadamard spaces (metric spaces with nonpositive curvature), the Fréchet mean is always unique \cite[Proposition 4.3]{sturm03}. This is not true in general. On the circle, a uniform distribution on two antipodal points has two Fréchet means. For a deeper analysis of Fréchet means on the circle, see \cite{hotz15}. Similarly, Fréchet means on many positively curved spaces like (hyper-)spheres may not be unique. For the metric space $\mc Q = \R$ with $d(q,p) =\sqrt{\abs{q-p}}$, the Fréchet mean set is the set of medians, which may also be non-unique. These examples underline the importance of considering sets of minimizers in a general theory instead of assuming uniqueness.

The notion of \textit{Fréchet mean} can be generalized to cases where the cost function to be minimized is not a squared metric, e.g.\ \cite{huckemann11}. We will not explicitly write down measurablity conditions, but silently demand that all spaces have the necessary measurable structure and all functions are measurable when necessary. Let $(\mc Q, d)$ be a metric space and $\mc Y$ be a set. Let $\mf c \colon \mc Y \times \mc Q \to \R$ be a function. Let $Y$ be a random variable with values in $\mc Y$. Let $M := \argmin_{q\in\mc Q} \Ex{\mf c(Y, q)}$ assuming the expectations exist. 
In this context, $\mf c$ is called \textit{cost function}, $\mc Y$ is called \textit{data space}, $\mc Q$ is called \textit{descriptor space}, $q\mapsto\Ex{\mf c(Y, q)}$ is called \textit{objective function} (or \textit{Fréchet function}), and $M$ is called \textit{generalized Fréchet mean set} or \textit{$\mf c$-Fréchet mean set}. 

This general scenario contains the setting of general M-estimation. It includes many important statistical frameworks like maximum likelihood estimation, where $\mc Q = \Theta$ parameterizes a family of densities $(f_\vartheta)_{\vartheta\in\Theta}$ on $\mc Y = \R^p$ and $\mf c(x,\vartheta) = -\log f_{\vartheta}(x)$, or linear regression, where $\mc Q = \R^{s+1}$, $\mc Y = (\{1\}\times\R^s) \times \R$, $\mf c((x,y),\beta) = (y - \beta\tr x)^2$. It also includes nonstandard settings, e.g.\ \cite{huckemann11}, where geodesics in $\mc Q$ are fitted to points in $\mc Y$.

Fix an arbitrary element $o\in\mc Q$. We will use cost functions $\mf c(y, q) = H(\ol yq)-H(\ol yo)$, where $H(x) = \int_0^x h(t) \dl t$ for a non-decreasing function $h$, and $\mf c(y, q) = \ol yq^\alpha-\ol yo^\alpha$ with $\alpha>0$. In both cases the set of minimizers does not depend on $o$. We call the minimizers of the former cost function \textit{$H$-Fréchet means}. In the latter case, we call the minimizers \textit{power Fréchet means} or \textit{$\alpha$-Fréchet means}. We can interpret the different exponents $\alpha = 2$, $\alpha =1$, $\alpha \to 0$, $\alpha \to \infty$ as mean, median, mode, and circumcenter (or mid-range), respectively, see \cite{macqueen67}. The minimizers for $\alpha=1$ are sometimes called \textit{Fréchet median}, e.g.\ \cite{arnaudon13}. If $\mc Q$ is a Banach space, then they are called \textit{geometric} or \textit{spatial median}, e.g.\ \cite{kemperman87}.
$H$-Fréchet means serve as a generalization of $\alpha$-Fréchet means for $\alpha>1$ as well as an intermediate result for proving strong laws of large numbers for $\alpha$-Fréchet mean sets with $\alpha\in(0,1]$.

For a function $f\colon\mc Q\to\R$ and $\epsilon\geq0$, define 
\begin{equation*}
	\epsilon\text{-}\argmin_{q\in\mc Q} f(q) := \Set{q \in \mc Q \given  f(q) \leq \epsilon +\inf_{q\in\mc Q} f(q) }
	\eqfs
\end{equation*}
Let $Y_1, \dots, Y_n$ be independent random variables with the same distribution as $Y$. Choose $(\epsilon_n)_{n\in\N}\subset [0,\infty)$ with $\epsilon_n \xrightarrow{n\to\infty}0$. Let $M_n := \epsilon_n\text{-}\argmin_{q\in\mc Q} \frac1n \sum_{i=1}^n \mf c(Y_i, q)$. Our goal is to show almost sure convergence of elements in $M_n$ to elements in $M$. 
\begin{remark}
	Considering sets of elements that minimize the objective only up to $\epsilon_n$ makes the results more relevant to applications in which Fréchet mean sets are approximated numerically.		
	Furthermore, it may allow us to find more elements of $M$ in the limit of $M_n$ than for $\epsilon_n=0$, as discussed in \autoref{rem:epsilon_argmin} below and appendix \ref{sec:median}.
\end{remark}
There are different possibilities of how a convergence of sets $M_n$ to a set $M$ can be described.
\begin{definition}
	Let $(\mc Q, d)$ be a metric space.
	\begin{enumerate}[label=(\roman*)]
		\item 
		Let $(B_n)_{n\in\N}$ with $B_n \subset \mc Q$ for all $n\in\N$. Then the \emph{outer limit} of $(B_n)_{n\in\N}$ is
		\begin{equation*}
			\outerlim_{n\to\infty} B_n := \bigcap_{n\in\N} \overline{\bigcup_{k\geq n} B_k}\eqcm
		\end{equation*}
		where $\overline B$ denotes the closure of the set $B$.
		\item 
		The \emph{one-sided Hausdorff distance} between $B, B^\prime \subset \MS$ is
		\begin{equation*}
			d_\subset(B,B^\prime) := \sup_{x\in B} \inf_{x^\prime\in B^\prime} d(x,x^\prime)\eqfs
		\end{equation*}
		\item 
		The \emph{Hausdorff distance} between $B, B^\prime \subset \MS$ is
		\begin{equation*}
			d_\Hausdorff(B, B^\prime) 
			:= 
			\max(d_\subset(B,B^\prime), d_\subset(B^\prime,B))
			\eqfs
		\end{equation*}
	\end{enumerate}
\end{definition}
\begin{remark}
\mbox{ }
\begin{enumerate}[label=(\roman*)]
	\item 
		The outer limit is the set of all points of accumulation of all sequences $(x_n)_{n\in\N}$ with $x_n\in B_n$. We may write $\outerlim_{n\to\infty} B_n = \Set{q\in\mc Q \given \liminf_{n\to\infty} d(B_n, q) = 0}$, where $d(B, q) := \inf_{p\in B} d(p, q)$ for a subset $B\subset\mc Q$. The \textit{inner limit} is dual to the outer limit. It is defined as $\innerlim_{n\to\infty} B_n := \Set{q\in\mc Q \given \limsup_{n\to\infty} d(B_n, q) = 0}$. Clearly, $\innerlim_{n\to\infty} B_n \subset \outerlim_{n\to\infty} B_n$. Thus, results of the form $\outerlim_{n\to\infty} B_n \subset B$, which we show below, are stronger than $\innerlim_{n\to\infty} B_n \subset B$.
	\item 
		It holds $d_\subset(B,B^\prime)=0$ if and only if $B \subset \overline{B^\prime}$, but $d_\Hausdorff(B,B^\prime)=0$ if and only if $\overline B = \overline B^ \prime$. The function $d_\Hausdorff$ is a metric on the set of closed and bounded subsets of $\MS$.
	\item 
		Elements from a sequence of sets might have sub-sequences that have no point of accumulation and are bounded away from the outer limit of the sequence of sets. That cannot happen with the one-sided Hausdorff limit. Here, every sub-sequence is eventually arbitrarily close to the limiting set. As an example, the outer limit of the sequence of sets $\{0,n\}$, $n\in\N$ on the Euclidean real line is $\{0\}$, but $d_\subset(\{0,n\},\{0\})\xrightarrow{n\to\infty}\infty$. Aside from an element with diverging distance ($n$ in the example), another cause for the two limits to not align may be non-compactness of bounded sets: Consider the space $\ell^2$ of all sequences $(x_k)_{k\in\N}\subset\R$ with $\sum_{k=1}^\infty x_k^2 < \infty$ with distance $d((x_k)_{k\in\N}, (y_k)_{k\in\N}) = (\sum_{k=1}^\infty (x_k-y_k)^2)^\frac12$. Let $\underline{0}\in\ell^2$ be the sequence with all entries equal to 0. Let $e^n := (e^n_k)_{k\in\N}\in\ell^2$ with $e_n^n=1$ and $e^n_k=0$ for all $k\neq n$. Then $d_\subset(\{\underline 0, e_n\}, \{\underline 0\}) = 1$ for all $n\in\N$, but $\outerlim_{n\to\infty} \{\underline 0, e_n\} = \{\underline 0\}$.
\end{enumerate}
\end{remark}
We will state conditions so that $\outerlim_{n\to\infty} M_n \subset M$ almost surely or $d_\subset(M_n, M) \xrightarrow{n\to\infty}_{\ms{a.s.}} 0$, where the index $\ms{a.s.}$ indicates almost sure convergence. It is not easily possible to show $d_\Hausdorff(M_n, M) \xrightarrow{n\to\infty}_{\ms{a.s.}} 0$ if $M$ is not a singleton, as discussed in \autoref{rem:epsilon_argmin} below and appendix \ref{sec:median}. These limit theorems may be called strong laws of large numbers of the Fréchet mean set or (strong) consistency of the empirical Fréchet mean set. Notably, in \cite{evans20} the connection to convergence in the sense of topology is made:
If the set of closed subsets of $\MS$ is equipped with the \textit{Kuratowski upper topology}, a sequence of closed subsets $(B_n)_{n\in\N}$ converges to a closed subset $B$ if and only if $\outerlim_{n\to\infty} B_n \subset B$. 
If the set of nonempty compact subsets of $\MS$ is equipped with the \textit{Hausdorff upper topology}, a sequence of nonempty compact subsets $(B_n)_{n\in\N}$ converges to a nonempty compact subset $B$ if and only if $d_{\subset}(B_n, B)\xrightarrow{n\to\infty} 0$. 

\cite{ziezold77} shows a strong law in outer limit for Fréchet mean sets with a second moment condition.
\cite{sverdrup81} shows a strong law in outer limit for power Fréchet mean sets in compact spaces.
\cite{bhattacharya03} shows almost sure convergence of Fréchet mean sets in one-sided Hausdorff distance with a second moment condition.
The independent parallel work \cite{evans20} shows strong laws in outer limit and one-sided Hausdorff distance for $\alpha$-Fréchet mean sets requiring $\Ex{\ol Yo^{\alpha}} < \infty$, which is a second moment condition for the Fréchet mean.
In contrast, we show strong laws of large numbers for power Fréchet mean sets in outer limit and in one-sided Hausdorff distance with less moment assumptions: For power $\alpha>1$, we require $\Ex{\ol Yo^{\alpha-1}} < \infty$, and for $\alpha\in (0,1]$ no moment assumption is made, see \autoref{cor:cons_da} and \autoref{cor:median}. Thus, $\alpha$-Fréchet means may be of interest in robust statistics.
\cite{huckemann11} shows almost sure convergence in one-side Hausdorff distance as well as in outer limit for generalized Fréchet means. Our results for $\mf c$-Fréchet means require slightly less strict assumptions, see \autoref{thm:epi} and \autoref{thm:consistency}, which make them applicable in a larger class of settings and allows us to derive our results for $H$- and $\alpha$-Fréchet means with minimal moment assumptions.
Results in \cite{artstein95, korf01, choirat03} imply strong laws and ergodic theorems in outer limit for generalized Fréchet means. We recite parts of these results to state \autoref{thm:epi}.
Furthermore, we show strong laws of large numbers for $H$-Fréchet means sets in outer limit, \autoref{cor:nondec:epi}, and one-sided Hausdorff distance, \autoref{cor:nondec:onehaus}.
When $M$ is singleton a quantitative version (rates of convergence) of the results presented in this article is given in \cite{schoetz19}.

Before we consider the probabilistic setting, we present theory on convergence of minimizing sets for deterministic functions in section \ref{sec:det}, where we partially follow \cite{rockafellar98}. Thereafter, we derive strong laws of large numbers for $\mf c$-Fréchet mean sets in section \ref{sec:gen}, for $H$-Fréchet mean sets in section \ref{sec:nondec}, and for $\alpha$-Fréchet mean sets in section \ref{sec:power}. 
Appendix \ref{sec:median} uses the median as a simple example to illustrate some peculiarities when dealing with sets of Fréchet means.
All strong laws in the main part of this article build upon \cite[Theorem 1.1]{korf01} -- a deep convergence result for functions $\MS \to \R$. In appendix \ref{sec:alt}, we show a different route to a strong law in one-side Hausdorff distance. This is illustrative, but requires slightly stricter assumptions. 
In appendix \ref{sec:aux} some auxiliary results are stated and proven.
 \section{Convergence of Minimizer Sets of Deterministic Functions}\label{sec:det}
Let $(\mc Q, d)$ be a metric space. The diameter of a set $B\subset \mc Q$ is defined as $\diam(B) = \sup_{q,p\in B} d(q,p)$. The following notion of convergence of functions will be useful to infer a convergence results of their minimizers.
\begin{definition}
	Let $f,f_n \colon \MS\to\R$, $n\in\N$.
	The sequence $(f_n)_{n\in\N}$ \emph{epi-converges} to $f$ \emph{at} $x\in\MS$ if and only if
	\begin{align*}
	\forall (x_n)_{n\in\N} \subset \MS, x_n \to x \colon& \liminf_{n\to\infty} f_n(x_n) \geq f(x)\qquad\text{and}\\
	\exists (y_n)_{n\in\N} \subset \MS, y_n \to x \colon& \limsup_{n\to\infty} f_n(y_n) \leq f(x)
	\eqfs
	\end{align*}
	The sequence $(f_n)_{n\in\N}$ \emph{epi-converges} to $f$ if and only if it epi-converges at all $x\in\MS$. We then write $f_n\epiconv{n}f$.
\end{definition}
We introduce some short notation. Let $f \colon \MS \to \R$ and $\epsilon\geq 0$. Denote $\inf f = \inf_{x\in\MS}f(x)$, $\argmin f = \Set{x \in\MS \given f(x) = \inf f}$, $\epsilon\text{-}\argmin f = \Set{x\in\MS \given  f(x) \leq \epsilon +\inf f }$.
Let $\delta > 0$ and $x_0\in\mc Q$.
Denote $\ball_\delta(x_0) = \Set{x\in\MS \given  d(x, x_0) < \delta}$.
Furthermore, $f$ is called \emph{lower semi-continuous} if and only if  $\liminf_{x\to x_0} f(x) \geq f(x_0)$ for all $x_0\in\mc Q$. 


To state convergence results for minimizing sets of deterministic functions, we need one final definition.
\begin{definition}
	A sequence $(B_n)_{n\in\N}$ of sets $B_n \subset \MS$ is called \emph{eventually precompact} if and only if there is $n\in\N$ such that the set 
$\bigcup_{k=n}^\infty B_k$ is precompact, i.e. its closure is compact.
\end{definition}
The first theorem of this section relates epi-convergence of functions to convergence of their sets of minimizers in outer limit.
\begin{theorem}\label{thm:convOfMini}
	Let $f,f_n \colon \MS\to\R$. Let $(\epsilon_n)_{n\in\N}\subset [0,\infty)$ with $\epsilon_n \xrightarrow{n\to\infty}0$. 
		Assume $f_n\epiconv{n}f$.
		Then
		\begin{equation*}
			\outerlim_{n\to\infty}\, \epsilon_n\text{-}\argmin f_n \subset \argmin f
		\end{equation*}
		and
		\begin{equation*}
			\limsup_{n\to\infty} \inf f_n \leq \inf f
			\eqfs
		\end{equation*}
\end{theorem}
Large parts of this theorem can be found e.g., in \cite[chapter 7]{rockafellar98}. To make this article more self-contained, we give a proof here.
\begin{proof}
		Let $x \in \outerlim_{n\to\infty}\, \epsilon_n\text{-}\argmin f_n$. Then there is a sequence $x_n \in \epsilon_{n}\text{-}\argmin f_{n}$ with a subsequence converging to $x$, i.e., $x_{n_i} \xrightarrow{i\to\infty} x$, where $n_i\xrightarrow{i\to\infty}\infty$.
		Let $y\in\MS$ be arbitrary. As $f_n\epiconv{n}f$, there is a sequence $(y_n)_{n\in\N}\subset\MS$ with $y_n\xrightarrow{n\to\infty}y$ and $\limsup_{n\to\infty} f_n(y_n) \leq f(y)$. 
		It holds $f_{n_i}(x_{n_i}) \leq \epsilon_{n_i} + \inf f_{n_i} \leq \epsilon_{n_i} + f_{n_i}(y_{n_i})$. Thus, by the definition of epi-convergence and $\epsilon_n \xrightarrow{n\to\infty} 0$, we obtain
		\begin{equation*}
		f(x) 
		\leq 
		\liminf_{i\to\infty} f_{n_i}(x_{n_i}) 
		\leq 
		\liminf_{i\to\infty} \br{\epsilon_{n_i} + f_{n_i}(y_{n_i})}
		\leq 
		\limsup_{i\to\infty} f_{n_i}(y_{n_i}) 
		\leq 
		f(y)\eqfs
		\end{equation*}		
		Thus, $x \in \argmin f$.
		Next, we turn to the inequality of the infima. For $\epsilon > 0$ choose an arbitrary $x\in\epsilon\text{-}\argmin f$. There is a sequence $(y_n)_{n\in\N}\subset\MS$ with $y_n\xrightarrow{n\to\infty} x$ and $f_n(y_n)\xrightarrow{n\to\infty} f(x)$. Thus,
		\begin{equation*}
			\limsup_{n\to\infty} \inf f_n \leq \limsup_{n\to\infty} f_n(y_n) \leq \inf f + \epsilon
			\eqfs
		\end{equation*}
\end{proof}
It is illustrative to compare this result with \autoref{thm:alt:convOfMini}, which shows that a stronger notion of convergence for functions -- convergences uniformly on bounded sets -- yields convergence of sets of minimizers in one-sided Hausdorff distance, which is a stronger notion of convergence of sets as the next theorem shows.
\begin{theorem}\label{thm:outerVsHaus}
	Let $(B_n)_{n\in\N}$ with $B_n \subset \mc Q$ for all $n\in\N$. Let $B \subset \MS$.
	\begin{enumerate}[label=(\roman*)]
		\item If $d_{\subset}(B_n, B) \xrightarrow{n\to\infty} 0$ then $\outerlim_{n\to\infty}\, B_n \subset \overline{B}$.
		\item Assume $(B_n)_{n\in\N}$ is eventually precompact. If $\outerlim_{n\to\infty}\, B_n \subset \overline{B}$ then $d_{\subset}(B_n, B) \xrightarrow{n\to\infty} 0$.
	\end{enumerate}	
\end{theorem}
\begin{proof}
	\mbox{ }
	\begin{enumerate}[label=(\roman*)]
\item 
	Assume $d_{\subset}(B_n, B) \xrightarrow{n\to\infty} 0$. Let $x_\infty \in \outerlim_{n\to\infty}\, B_n$, i.e., there is a sequence $(x_{n_k})_{k\in\N}\subset \MS$ with $n_1 < n_2 < \dots$ and $x_{n_k}\in B_{n_k}$ such that $x_{n_k} \xrightarrow{k\to\infty} x_\infty$. Thus, 
	\begin{equation*}
		\inf_{x\in B} d(x_\infty, x) \leq d(x_\infty, x_{n_k}) + \inf_{x\in B} d(x_{n_k}, x)  \xrightarrow{k\to\infty} 0\eqfs
	\end{equation*}
	This shows $\inf_{x\in B} d(x_\infty, x) = 0$. Hence,
	\begin{equation*}
		\outerlim_{n\to\infty}\,B_n \subset \Set{x_\infty \in \MS \given \inf_{x\in B} d(x_\infty, x) = 0} = \overline{B}\eqfs
	\end{equation*}
\item	
	Assume  $\outerlim_{n\to\infty}\, B_n \subset \overline{B}$. Further assume the existence of $\epsilon>0$ and a sequence $(x_{n_k})_{k\in\N}\subset \MS$ with $n_1 < n_2 < \dots$ and $x_{n_k}\in B_{n_k}$ such that $\inf_{x\in B} d(x_{n_k}, x) \geq \epsilon$. As $(B_{n_k})_{k\in\N}$ is eventually precompact, the sequence $(x_{n_k})_{k\in\N}$ has an accumulation point $x_\infty$ in $\overline{\bigcup_{k \geq k_0} B_{n_k}}$ for some $k_0\in\N$ with $\inf_{x\in B} d(x_{\infty}, x) \geq \epsilon$. In particular, $x_{\infty} \not\in \overline{B}$, which contradicts the first assumption in the proof. Thus, a sequence $(x_{n_k})_{k\in\N}$ with these properties cannot exist, which implies $d_{\subset}(B_n, B) \xrightarrow{n\to\infty} 0$.
	\end{enumerate}
\end{proof}
Note that the argument for the second part is essentially the same as in \cite[proof of Theorem A.4]{huckemann11}.
\begin{remark}\label{rem:epsilon_argmin}
Together \autoref{thm:convOfMini} and \autoref{thm:outerVsHaus} may yield convergence of minimizers in one-sided Hausdorff distance. But even if $d_{\subset}(\epsilon_n\text{-}\argmin f_n, \argmin f) \xrightarrow{n\to\infty} 0$, $d_\Haus(\epsilon_n\text{-}\argmin f_n, \argmin f)$ does not necessarily vanish unless $\argmin f$ is a singleton.
Similarly, for an arbitrary sequence $\epsilon_n \xrightarrow{n\to\infty}0$, the outer limit of $\epsilon_n\text{-}\argmin f_n$ may be a strict subset of $\argmin f$.
But according to \cite[Theorem 7.31 (c)]{rockafellar98}, there exists a sequence $(\epsilon_n)_{n\in\N}$ with $\epsilon_n\xrightarrow{n\to\infty}0$ slow enough such that $\outerlim_{n\to\infty}\epsilon_n\text{-}\argmin f_n = \argmin f$. An explicit example of this phenomenon is presented in appendix \ref{sec:median}.
\end{remark}
 \section{Strong Laws for $\mf c$-Fréchet Mean Sets}\label{sec:gen}
Let $(\mc Q, d)$ be a metric space, the descriptor space. Let $\mc Y$ be a set, the data space. Let $\mf c\colon \mc Y \times\mc Q \to \R$ be a function, the cost function. Let $(\Omega, \Sigma, \Pr)$ be a probability space that is silently underlying all random variables in this section. Let $Y$ be a random variable with values in $\mc Y$. Denote the $\mf c$-Fréchet mean set of $Y$ as $M = \argmin_{q\in\mc Q} \Ex{\mf c(Y, q)}$.
Let $Y_1, \dots, Y_n$ be independent random variables with the same distribution as $Y$.
Choose $(\epsilon_n)_{n\in\N}\subset [0,\infty)$ with $\epsilon_n \xrightarrow{n\to\infty}0$. Set $M_n = \epsilon_n\text{-}\argmin_{q\in\mc Q} \frac1n \sum_{i=1}^n \mf c(Y_i, q)$.
\begin{assumptions}\mbox{ }
\begin{itemize}
\item 
	\textsc{Polish}: 
	$(\mc Q, d)$ is separable and complete.
\item 
	\textsc{LowerSemiContinuity}: 
	$q \mapsto \mf c(y, q)$ is lower semi-continuous.
\item 
	\textsc{Integrable}: 
	$\Ex{\abs{\mf c(Y, q)}} < \infty$ for all $q\in\mc Q$.
\item 
	\textsc{IntegrableInf}: 
	$\Ex{\inf_{q \in \mc Q}\mf c(Y, q)} > - \infty$.
\end{itemize}
\end{assumptions}
\begin{theorem}\label{thm:epi}
	Assume \textsc{Polish}, \textsc{LowerSemiContinuity}, \textsc{Integrable}, and \textsc{IntegrableInf}.
	Then, almost surely,
	\begin{equation*}
		\outerlim_{n\to\infty}\, M_n \subset M
		\eqfs
	\end{equation*}
\end{theorem}
\begin{proof}
Define $F(q) = \Ex{\mf c(Y, q)}$, $F_n(q) = \frac1n\sum_{i=1}^n \mf c(Y_i, q)$. By \textsc{Integrable}, $F(q)<\infty$.
\cite[Theorem 1.1]{korf01} states that $F_n \xrightarrow{n\to\infty}_{\ms{epi}} F$ almost surely if \textsc{Polish}, \textsc{LowerSemiContinuity}, and \textsc{IntegrableInf} are true. \autoref{thm:convOfMini} then implies $\outerlim_{n\to\infty}\, M_n \subset M$ almost surely.
\end{proof}
\begin{assumptions}\mbox{ }
\begin{itemize}
	\item \textsc{HeineBorel}: Every closed bounded set in $\mc Q$ is compact.
	\item \textsc{SampleHeineBorel}: Almost surely following is true: There is $N_0\in\N$ such that every closed and bounded subset of $\bigcup_{n\geq N_0} M_n$ is compact.
	\item \textsc{UpperBound}: $\Ex{\sup_{q\in B} \abs{\mf c(Y,q)}} < \infty$ for all bounded sets $B\subset\mc Q$.
	\item \textsc{LowerBound}: There are $o\in\mc Q$, $\psi^+, \psi^- \colon [0,\infty) \to [0,\infty)$, $\mf a^+, \mf a^- \in(0,\infty)$, and $\sigma(Y_1, \dots, Y_n)$-measurable random variables $\mf a^+_n, \mf a^-_n \in[0,\infty)$ such that 
	\begin{align*}
		\mf a^+ \psi^+(\ol qo) - \mf a^- \psi^-(\ol qo) &\leq \Ex{\mf c(Y, q)}
		\eqcm\\
		\mf a^+_n \psi^+(\ol qo) - \mf a^-_n \psi^-(\ol qo) &\leq \frac1n \sum_{i=1}^n \mf c(Y_i, q)
	\end{align*}
	for all $q\in\mc Q$. Furthermore, $\mf a^+_n\xrightarrow{n\to\infty}_{\ms{a.s.}}\mf a^+$ and $\mf a^-_n\xrightarrow{n\to\infty}_{\ms{a.s.}}\mf a^-$. Lastly, $\psi^+(\delta)/\max(1, \psi^-(\delta))\xrightarrow{\delta\to\infty}\infty$.\footnote{A previous version omitted the $\max(1,\cdot)$, which was an error. I became aware of this mistake through work by Jaesung Park and Sungkyu Jung, whom I thank for their careful analysis.}
	\end{itemize}
\end{assumptions}
\begin{remark}\label{rem:on_ass1}\mbox{ }
\begin{itemize}
\item 
	Following implications hold:
	\begin{align*}
	\textsc{HeineBorel} &\Rightarrow \textsc{Polish}\eqcm\\
	\textsc{HeineBorel} &\Rightarrow \textsc{SampleHeineBorel}\eqcm\\
	\textsc{UpperBound} &\Rightarrow \textsc{Integrable}\eqfs
	\end{align*}
\item On \textsc{HeineBorel}:
	A space enjoying this property is also called \textit{boundedly compact} or \textit{proper} metric space.
	The Euclidean spaces $\R^s$, finite dimensional Riemannian manifolds, as well as $\mc C^\infty(U)$ for open subsets $U \subset \R^s$  fulfill \textsc{Heine--Borel} \cite[section 8.4.7]{edwards95}.
	See \cite{williamson87} for a construction of further spaces where \textsc{Heine--Borel} is true. 
\item On \textsc{SampleHeineBorel} and infinite dimension:
	If $M_n=\{m_n\}$ and $M=\{m\}$ are singleton sets and $m_n \xrightarrow{n\to\infty} m$ almost surely, then \textsc{SampleHeineBorel} holds. It is less strict than \textsc{HeineBorel}: In separable Hilbert spaces of infinite dimension \textsc{HeineBorel} does not hold. But with the metric $d$ induced by the inner product and $\mf c=d^2$, $\mf c$-Fréchet means are unique and equal to the usual notion of mean. Furthermore, strong laws of large numbers in Hilbert spaces are well-known, see e.g.\ \cite{kawabe86}. Thus, \textsc{SampleHeineBorel} is true. Let it be noted that proving \textsc{SampleHeineBorel} in a space where \textsc{HeineBorel} is false may be of similar difficulty as showing convergence of Fréchet means directly. In the case of infinite dimensional Banach spaces, results on strong laws of large numbers for a different notion of mean -- the Bochner integral -- are well established, see e.g.\ \cite{hoffmann76}.
\item On \textsc{LowerBound}:
	We illustrate this condition in the linear regression setting with $\mc Q = \R^{s+1}$, $\mc Y = (\{1\}\times\R^s) \times \R$, $\mf c((x,y),\beta) = (y - \beta\tr x)^2 - y^2 = - 2 \beta\tr x y + \beta\tr xx\tr \beta$. Let $(X, Y)$ be random variables with values in $\mc Y$. Let $(X_1, Y_1),\dots, (X_n, Y_n)$ be independent with the same distribution as $(X,Y)$. We can set $o = 0\in\R^{s+1}$, $\mf a^+ = \lambda_{\ms{min}}(\Ex{XX\tr})$, where $\lambda_{\ms{min}}$ denotes the smallest eigenvalue, $\mf a^- = 2 \normof{\Ex{XY}}$, $\mf a^+_n = \lambda_{\ms{min}}(\frac1n \sum_{i=1}^n X_iX_i\tr)$, $\mf a^-_n = 2 \normof{\frac1n\sum_{i=1}^n X_iY_i}$, $\psi^+(\delta) = \delta^2$ and $\psi^-(\delta) = \delta$. If $\lambda_{\ms{min}}(\Ex{XX\tr}) > 0$, the largest eigenvalue $\lambda_{\ms{max}}(\Ex{XX\tr}) < \infty$, and $\Ex{\normof{XY}} < \infty$, all conditions are fulfilled.
	
	For a further application of \textsc{LowerBound}, see the proof of \autoref{cor:nondec:onehaus} in the next section.
\end{itemize}
\end{remark}
\begin{theorem}\label{thm:consistency}
	Assume \textsc{Polish}, \textsc{LowerSemiContinuity}, \textsc{IntegrableInf},  \textsc{SampleHeineBorel}, \textsc{UpperBound}, and \textsc{LowerBound}. Then
	\begin{equation*}
		d_\subset(M_n, M) \xrightarrow{n\to\infty}_{\ms{a.s.}} 0
		\eqfs
	\end{equation*}
\end{theorem}
\begin{proof}
The proof consists of following steps:
\begin{enumerate}
\item Apply \autoref{thm:epi}.
\item Reduction to a bounded set.
\item Show that $M_n$ is eventually precompact almost surely.
\item Apply \autoref{thm:outerVsHaus}.
\end{enumerate}
\noindent
\underline{\smash{Step 1.}}
\textsc{Polish}, \textsc{LowerSemiContinuity}, and \textsc{IntegrableInf} are assumptions. \textsc{UpperBound} implies \textsc{Integrable}. Thus, \autoref{thm:epi} yields $\outerlim_{n\to\infty}\, M_n \subset M$ almost surely.

\noindent
\underline{\smash{Step 2.}}
Define $F(q) = \Ex{\mf c(Y, q)}$, $F_n(q) = \frac1n\sum_{i=1}^n \mf c(Y_i, q)$. 
We want to show that there is a bounded set $B_1 \subset \mc Q$ such that $F(q) \geq F(m) + 1$ and $F_n(q) \geq F_n(m) + 1$ for all $q\in \mc Q \setminus B_1$ and $m \in M$. If $\mc Q$ is bounded, we can take $B_1 = \mc Q$. Assume $\mc Q$ is not bounded. 

Let $m \in M$. By \textsc{UpperBound}, $F(m) < \infty$. Let $o\in\mc Q$ from \textsc{LowerBound}.
Due to \textsc{LowerBound}, $F(q) \geq \mf a^+ \psi^+(\delta) - \mf a^- \psi^-(\delta) \geq  F(m) + 2$ for all $q\in\mc Q\setminus \ball_\delta(o)$ and $\delta$ large enough. This holds for all $m\in M$ as $F(m)$ does not change with $m$. We set $B_1 = \ball_\delta(o)$. For $F_n$, it holds $F_n(m) \xrightarrow{n\to\infty}_{\ms{a.s.}} F(m)$ and $\inf_{q\in\mc Q \setminus B_1} F_n(q) \geq \mf a^+_n \psi^+(\delta) - \mf a^-_n \psi^-(\delta)$ with $\mf a^+_n \xrightarrow{n\to\infty}_{\ms{a.s.}} \mf a^+$ and $\mf a^-_n \xrightarrow{n\to\infty}_{\ms{a.s.}} \mf a^-$.
Thus, there is a random variable $N_1$ such that almost surely $F_n(q) \geq F_n(m) + 1$ for all $n \geq N_1$, $q\in \mc Q \setminus B_1$, and $m\in M$.

\noindent
\underline{\smash{Step 3.}}
Take $N_0$ from \textsc{SampleHeineBorel}. Choose $N_2\geq \max(N_0, N_1)$ such that $\epsilon_n < 1$ for all $n \geq N_2$. Then $M_n \subset B_1$ for all $n\geq N_2$. Thus, $\bigcup_{n\geq N_2}M_n$ is bounded and -- due to \textsc{SampleHeineBorel}, \textsc{Polish}, and \autoref{lmm:precompact} -- precompact almost surely. 

\noindent
\underline{\smash{Step 4.}}
Finally, step 1 and 3 together with \autoref{thm:outerVsHaus} yield $d_\subset(M_n, M) \xrightarrow{n\to\infty}_{\ms{a.s.}} 0$.
\end{proof}
 \section{Strong Laws for $H$-Fréchet Mean Sets}\label{sec:nondec}
Let $(\mc Q, d)$ be a metric space. Let $(\Omega, \Sigma, \Pr)$ be a probability space that is silently underlying all random variables in this section. Let $Y$ be a random variable with values in $\mc Q$.
Let $h \colon [0,\infty) \to [0,\infty)$ be a non-decreasing function.
Define $H \colon [0,\infty) \to [0,\infty), x\mapsto \int_0^x h(t) \dl t$.
Fix an arbitrary element $o\in\mc Q$. Denote the $H$-Fréchet mean set of $Y$ as $M = \argmin_{q\in\mc Q} \Ex{H(\ol Yq) - H(\ol Yo)}$.
Let $Y_1, \dots, Y_n$ be independent random variables with the same distribution as $Y$.
Choose $(\epsilon_n)_{n\in\N}\subset [0,\infty)$ with $\epsilon_n \xrightarrow{n\to\infty}0$. Set $M_n = \epsilon_n\text{-}\argmin_{q\in\mc Q} \frac1n \sum_{i=1}^n (H(\ol {Y_i}q) - H(\ol {Y_i}o))$.
\begin{assumptions}\mbox{ }
	\begin{itemize}
	\item \textsc{InfiniteIncrease}:
		$h(x) \xrightarrow{x\to\infty} \infty$.
	\item \textsc{Additivity}:
		There is $b \in [1,\infty)$ such that $h(2x) \leq b h(x)$ for all $x \geq 0$.
	\item \textsc{$h$-Moment}: 
		$\Ex{h(\ol Yo)} < \infty$.
	\end{itemize}
\end{assumptions}
\begin{remark}\mbox{ }
\begin{itemize}
\item On \textsc{Additivity}:
	This implies $h(x+y)\leq b(h(x) + h(y))$ for all $x,y\geq 0$, see \autoref{lmm:nondec} (appendix).
	If $h$ is concave, \textsc{Additivity} holds with $b=2$ and we even have $h(x+y)\leq h(x) + h(y)$.
	This condition is not very restrictive, but it excludes functions that grow exponentially.
\end{itemize}
\end{remark}
\begin{corollary}\label{cor:nondec:epi}
	Assume \textsc{Polish}, \textsc{Additivity}, and \textsc{$h$-Moment}.
	Then, almost surely, 
	\begin{equation*}
		\outerlim_{n\to\infty}\, M_n \subset M
		\eqfs
	\end{equation*}
\end{corollary}
\begin{proof}
	We check the conditions of \autoref{thm:epi}. \textsc{Polish} is an assumption. 
	\textsc{LowerSemiContinuity} is fulfilled as $(q,p)\mapsto d(q,p)$ and $x \mapsto H(x)$ are continuous.
	For \textsc{Integrable}, we note that $H$ is non-decreasing and apply \autoref{lmm:nondec} (i),
	\begin{align*}
		\abs{H(\ol yq) - H(\ol yo)} 
		&\leq 
		\abs{\ol yq - \ol yo} h\brOf{\max(\ol yq, \ol yo)}
		\\&\leq 
		\ol qo\, h(\ol qo + \ol yo)
		\\&\leq 
		b\, \ol qo \br{h(\ol qo) + h(\ol yo)}
		\eqcm
	\end{align*}
	where the last inequality follows from \autoref{lmm:nondec} (ii) using \textsc{Additivity}.
	Thus, \textsc{$h$-Moment} implies \textsc{Integrable}.
	To show \textsc{IntegrableInf}, we note that $H$ is non-decreasing and apply \autoref{lmm:nondec} (iii),
	\begin{align*}
		H(\ol yq) - H(\ol yo)
		&\geq 
		H(\abs{\ol yo - \ol qo}) - H(\ol yo)
		\\&\geq 
		b^{-1} H( \ol qo) - 2  \,\ol qo\, h( \ol yo)
	\end{align*}
	due to \textsc{Additivity}. Furthermore, $H(\delta)  = \int_0^\delta h(x) \dl x \geq \frac12 \delta h(\frac12 \delta)$.
	With that, \textsc{$h$-Moment} implies \textsc{IntegrableInf}. 
	Thus, \autoref{thm:epi} can be applied.
\end{proof}
\begin{corollary}\label{cor:nondec:onehaus}
	Assume \textsc{SampleHeineBorel}, \textsc{Polish}, \textsc{Additivity}, \textsc{InfiniteIncrease}, and \textsc{$h$-Moment}.
	Then
	\begin{equation*}
		d_\subset(M_n, M) \xrightarrow{n\to\infty}_{\ms{a.s.}} 0
		\eqfs
	\end{equation*}
\end{corollary}
\begin{proof}
	We check the conditions of \autoref{thm:consistency}. \textsc{SampleHeineBorel} and \textsc{Polish} are assumptions of the corollary. \textsc{LowerSemiContinuity} and \textsc{IntegrableInf} are shown in the proof of \autoref{cor:nondec:epi}.
	Following that proof, we find, due to \textsc{Additivity},
	\begin{align*}
		\abs{H(\ol yq) - H(\ol yo)} &\leq b\, \ol qo \br{h(\ol qo) + h(\ol yo)}\eqcm
		\\
		H(\ol yq) - H(\ol yo) &\geq 	b^{-1} H( \ol qo) - 2 \,\ol qo\, h( \ol yo)\eqcm\\
		H(\delta)  &\geq \frac12 \delta h\brOf{\frac12 \delta}
		\eqfs
	\end{align*}
	The first inequality together with \textsc{$h$-Moment} implies \textsc{UpperBound}.
	For \textsc{LowerBound}, we use the second inequality. We set $\psi^+(\delta) = 	b^{-1} H(\delta)$, $\psi^-(\delta) = 2 \delta$, $\mf a^+=\mf a_n^+ = 1$, $\mf a^{-} = \Ex{h(\ol Yo)}$, and $\mf a^-_n = \frac1n \sum_{i=1}^n h(\ol {Y_i}o)$ with $\mf a^-_n\xrightarrow{n\to\infty}_{\ms{a.s.}}\mf a^-$ due to \textsc{$h$-Moment}. Because of the third inequality, $\psi^+(\delta)/\psi^-(\delta) \geq \frac14 b^{-1} h(\frac12 \delta) \xrightarrow{\delta\to\infty} \infty$ by \textsc{InfiniteIncrease}.
\end{proof}
 \section{Strong Laws for $\alpha$-Fréchet Mean Sets}\label{sec:power}
Let $(\mc Q, d)$ be a metric space. Let $(\Omega, \Sigma, \Pr)$ be a probability space that is silently underlying all random variables in this section. Let $Y$ be a random variable with values in $\mc Y$.
Let $\alpha > 0$. Fix an arbitrary element $o\in\mc Q$. Denote the $\alpha$-Fréchet mean set of $Y$ as $M = \argmin_{q\in\mc Q} \Ex{\ol Yq^\alpha - \ol Yo^\alpha}$.
Let $Y_1, \dots, Y_n$ be independent random variables with the same distribution as $Y$.
Choose $(\epsilon_n)_{n\in\N}\subset [0,\infty)$ with $\epsilon_n \xrightarrow{n\to\infty}0$. Set $M_n =\epsilon_n\text{-}\argmin_{q\in\mc Q} \frac1n \sum_{i=1}^n (\ol {Y_i}q^\alpha - \ol{Y_i}o^\alpha)$.
\begin{corollary}\label{cor:cons_da}
	Let $\alpha > 1$.
	Assume $\Ex{\ol Yo^{\alpha-1}} < \infty$ and \textsc{Polish}.	
	\begin{enumerate}[label=(\roman*)]
	\item Then $\outerlim_{n\to\infty}\, M_n \subset M$ almost surely.
	\item Additionally, assume \textsc{SampleHeineBorel}.
		Then $d_\subset(M_n, M) \xrightarrow{n\to\infty}_{\ms{a.s.}} 0$.
	\end{enumerate}	
\end{corollary}
\begin{proof}
	Set $h(x) = \alpha x^{\alpha-1}$. This function is non-decreasing, fulfills \textsc{Additivity} with $b = 2^{\alpha-1}$ and \textsc{InfiniteIncrease}.
	Due to $\Ex{\ol Yo^{\alpha-1}} < \infty$, \textsc{$h$-Moment} is fulfilled. Furthermore, $H(x) = x^\alpha$. Thus, \autoref{cor:nondec:onehaus} and \autoref{cor:nondec:epi} imply the claims.
\end{proof}
\begin{corollary}\label{cor:median}
	Let $\alpha \in (0,1]$. Assume \textsc{Polish}.
	\begin{enumerate}[label=(\roman*)]
	\item Then $\outerlim_{n\to\infty}\, M_n \subset M$ almost surely.
	\item Additionally, assume \textsc{SampleHeineBorel}.
		Then $d_\subset(M_n, M) \xrightarrow{n\to\infty}_{\ms{a.s.}} 0$.
	\end{enumerate}	
\end{corollary}
\begin{proof}
First, consider the case $\alpha=1$.
Apply \autoref{lmm:nondec:existence} (appendix) on $\ol Yo$ to obtain a function $h \colon [0,\infty) \to [0,\infty)$ which is strictly increasing, continuous, concave, fulfills \textsc{InfiniteIncrease}, and $\Ex{h(\ol Yo)} < \infty$. Concavity implies \textsc{Additivity} with $b=2$.
As its derivative is strictly increasing, $H(x) = \int_0^x h(t) \dl t$ is convex and strictly increasing. Thus, $H$ has an inverse $H^{-1}$ and $H^{-1}$ is concave. 
This implies that $d_H(q,p) = H^{-1}(\ol qp)$ is a metric. 

As $H^{-1}$ is concave, there are $u_0, u_1 \in[0,\infty)$ such that $H^{-1}(x) \leq u_0 + u_1 x$ for all $x\geq0$.
As $h$ is concave, there are $v_0, v_1\in[0,\infty)$ such that $h(u_0 + u_1 x) \leq v_0 + v_1 h(x)$ for all $x\geq 0$.
Thus, $\Ex{h(d_H(Y, o))} = \Ex{h(H^{-1} (\ol Yo))} \leq v_0 + v_1 \Ex{h(\ol Yo)} < \infty$.
Hence, \textsc{$h$-Moment} is true for the metric $d_H$.

Moreover, \textsc{Polish} and \textsc{HeineBorel}-type properties of $(\mc Q, d)$ are preserved in $(\mc Q, d_H)$, as $H^{-1}$ is strictly increasing, concave, and continuous, with $H^{-1}(0) = 0$ and $H^{-1}(\delta)\xrightarrow{\delta\to\infty}\infty$, and thus, the properties boundedness, compactness, separability, and completeness coincide for $d$ and $d_H$.
Applying \autoref{cor:nondec:epi} and \autoref{cor:nondec:onehaus} on the minimizers of $\Ex{H(d_H(Y, q)) - H(d_H(Y, o))}  = \Ex{\ol  Yq - \ol Yo}$ now yields the claims for $\alpha=1$.

For $\alpha \in (0,1)$ just note, that $\tilde d(q,p) = d(q,p)^\alpha$ is a metric, which preserves \textsc{Polish} and \textsc{HeineBorel}-type properties, and apply the result for $\alpha = 1$ on $\tilde d$.
\end{proof}
\begin{remark}
	 For convergence, we need the $(\alpha-1)$-moment to be finite in the case of $\alpha \geq 1$. \cite[Corollary 5]{schoetz19} shows that in metric spaces with nonnegative curvature the typical parametric rate of convergence $n^{-\frac12}$ is obtained for $\alpha$-Fréchet means assuming the $2(\alpha-1)$-moment to be finite in the case of $\alpha\in[1,2]$ under some further conditions.
\end{remark}
 \begin{appendix}
\section{Example: The Set of Medians}\label{sec:median}
Let $s\in\N$. Consider the metric space $(\R^s, d_1)$, where $d_1(q,p) = \normof{q-p}_1 = \sum_{j=1}^s \abs{q_j-p_j}$. The power Fréchet mean with $\alpha=1$ in this space is equivalent to the standard ($\alpha=2$) Fréchet mean in $(\R^s, d_1^{\frac12})$. For $s=1$ it is equal to the median.
Let $Y = (Y^1,\dots, Y^s)$ be a random vector in $\R^s$ such that $\Pr(Y^k=0)=\Pr(Y^k=1)=\frac12$ for $k=1,\dots, s$ and $Y^1,\dots, Y^s$ are independent. Let $Y_1, Y_2,\dots$ be independent and identically distributed copies of $Y$. Let $M = \argmin_{q\in\R^s}\Ex{d_1(Y, q)}$ be the Fréchet mean set of $Y$ and $M_n = \epsilon_n\text{-}\argmin_{q\in\R^s}\frac1n\sum_{i=1}^n d_1(Y_i, q)$ its sample version.

\subsection{No Convergence in Hausdorff Distance}
First consider the case $s=1$ and $\epsilon_n=0$.
As $s=1$,  $M = \argmin_{q\in\R} \Ex{\abs{Y-q}}$ is the median of $Y$, which is $M = [0,1]$ as $2\Ex{\abs{Y-q}} = \abs{1-q}+\abs{q}$ achieves its minimal value $1$ precisely for all $q\in[0,1]$.
Define $p_n := \frac1n \sum_{i=1}^n Y_i$. Then the empirical objective function is $F_n(q) := p_n\abs{1-q}+(1-p_n)\abs{q}$, i.e., the sample Fréchet mean set is $M_n = \argmin_{q\in\R} F_n(q)$.
If $n$ is odd, then either $M_n = \{0\}$ or $M_n = \{1\}$ holds. The same is true for an even value of $n$ except when $p_n=\frac12$, in which case $M_n=[0,1]$.
Thus, $d_{\subset}(M_n, M) = 0$, but $d_{\Hausdorff}(M_n, M)$ does not converge almost surely.

\subsection{The Outer Limit as a Strict Subset}
Next, we keep $\epsilon_n=0$, but consider the value of $\outerlim_{n\to\infty} M_n$ in a multi-dimensional setting, i.e., $s\in\N$, as this yields a potentially surprising result:
By \autoref{lmm:products}, $M$ is just the Cartesian product of the median sets in each dimension, i.e., $M = [0,1]^s$ (this is not to be confused with the geometric median, which is the Fréchet mean with respect to the square root of the Euclidean norm). Similarly, $M_n = \bigtimes_{k=1}^s M_n^k$ decomposes into the sample Fréchet mean sets $M_n^k$ of each dimension $k=1,\dots,s$. It holds $M_n^k = [0,1]$ if and only if the respective value of $p_n^k := \frac1n \sum_{i=1}^n Y_i^k$ is equal to $\frac12$, i.e., if and only if the symmetric simple random walk $S_n^k := \sum_{i=1}^n (2Y_i^k-1)$ hits $0$. Let $N = \#\Set{n\in\N\given S_n^1=\dots=S_n^s=0}$. Let $A\subset \R$ and $B_n\subset \R$ for all $n\in\N$. If $A\subset B_n$ for infinitely many $n$, then $A \subset \outerlim_{n\to\infty} B_n$. Thus, we want to know whether $N$ is finite or infinite. This is answered by \textit{Pólya's Recurrence Theorem} \cite{polya21}. It implies that for $s \in \{1, 2\}$, $N = \infty$ almost surely. Furthermore, if $s\geq 3$, then $N < \infty$ almost surely.
To find which points are not element of the outer limit of sample Fréchet mean sets, note following fact: For an open subset $A\subset \R$, if $A \subset \R \setminus B_n$ for all but finitely many $n$, then $A \subset \R \setminus \outerlim_{n\to\infty} B_n$.
We conclude, that in general a vector $x \in [0,1]^s$ is an element of $\outerlim_{n\to\infty} M_n$ if and only if at most two entries are not in $\{0,1\}$, i.e., almost surely
\begin{equation*}
	\outerlim_{n\to\infty} M_n = \Set{(x_1,\dots,x_s) \in [0,1]^s \given \#\{k\in\{1,\dots,s\} | x_k \in (0,1)\} \leq 2}
	\eqfs
\end{equation*}
Thus, $\outerlim_{n\to\infty} M_n = M$ for $s\in\{1,2\}$ and $\outerlim_{n\to\infty} M_n \subsetneq M$ for $s \geq 3$.

\subsection{Convergence in Hausdorff Distance}
Lastly, we use the setting $s=1$ and $\epsilon_n\in[0,\infty)$, where we want to find $\epsilon_n$ such that $[0,1] \subset M_n$.
At least one of $0$ and $1$ is a minimizer of $F_n(q)=p_n\abs{1-q}+(1-p_n)\abs{q}$ and the $F_n(q)$ is linear on $[0,1]$. Thus, $[0,1] \subset M_n$ if and only if  $\epsilon_n \geq |F_n(0)-F_n(1)|$. This is equivalent to 
 $|p_n-\frac12| < \frac12\epsilon_n$.
By Markov's inequality
\begin{equation*}
	\PrOf{\abs{p_n-\frac12} \geq \frac12\epsilon_n} \leq \frac{n^{-3} \Ex*{(Y-\frac12)^4}}{2^{-4}\epsilon_n^4}
	\eqfs
\end{equation*}
For $\epsilon_n = n^{-\frac14}$, we obtain
\begin{equation*}
	\sum_{n=1}^\infty \PrOf{\abs{p_n-\frac12} \geq \frac12\epsilon_n}  \leq \sum_{n=1}^\infty n^{-2} < \infty
	\eqfs
\end{equation*}
The Borel--Cantelli lemma implies that almost surely and for all $n$ large enough,  $|p_n-\frac12| < \frac12\epsilon_n$ and thus, $[0,1] \subset M_n$. 
Together with \autoref{cor:median}, we obtain $d_{\Hausdorff}(M_n, M) \xrightarrow{n\to\infty}_{\ms{a.s.}} 0$.
 \section{Alternative Route to One-Sided Hausdorff Convergence}\label{sec:alt}
In this section, we show an alternative proof of a strong law of large numbers for generalized Fréchet mean sets in one-sided Hausdorff distance. Although the final result, \autoref{thm:alt:consistency}, is weaker than \autoref{thm:consistency}, it is very illustrative to follow this line of proof:
In contrast to the arguments in the main part of the article, it does not rely on the powerful result \cite[Theorem 1.1]{korf01}, which seems to be rather complex to prove. Instead our reasoning here is simpler and more self-contained.
Furthermore, a comparison between convergence in outer limit and in one-sided Hausdorff distance seems more natural in view of the deterministic results \autoref{thm:convOfMini} and \autoref{thm:alt:convOfMini}, and the stochastic results \autoref{thm:epi} and \autoref{thm:alt:consistency}.
\subsection{Convergence of Minimizer Sets of Deterministic Functions}\label{ssec:alt:det}
Let $(\mc Q, d)$ be a metric space. For $A\subset \mc Q$ and $\delta>0$, denote $\ball_\delta(A) = \bigcup_{x\in A} \ball_\delta(x)$. 
\begin{definition} \mbox{ }
	\begin{enumerate}[label=(\roman*)]
	\item 
		Let $f,f_n \colon \MS\to\R$, $n\in\N$.
		The sequence $(f_n)_{n\in\N}$  \emph{converges} to $f$ \emph{uniformly on bounded sets} if and only if for every $B\subset \MS$ with $\diam(B)<\infty$,
		\begin{equation*}
		\lim_{n\to\infty} \sup_{x\in B}\abs{f_n(x) - f(x)} =0
		\eqfs
		\end{equation*}
		We then write $f_n\ubsconv{n}f$.
	\item 
		A sequence $(B_n)_{n\in\N}$ of sets $B_n \subset \MS$ is called \emph{eventually bounded} if and only if
		\begin{equation*}
		\limsup_{n\to\infty} \diam\brOf{\bigcup_{k=n}^\infty B_k} < \infty
		\eqfs
		\end{equation*}
	\item 
		A function $f$ has \emph{approachable minimizers} if and only if for all $\epsilon >0$ there is a $\delta>0$ such that
		$\delta\text{-}\argmin f \subset B_\epsilon(\argmin f)$.
	\end{enumerate}	
\end{definition}
The last definition directly implies that $d_\subset(\delta\text{-}\argmin f, \argmin f) \xrightarrow{\delta\to0} 0$ is equivalent to $f$ having approachable minimizers. Furthermore, if $f$ has approachable minimizers, then $\argmin f \neq \emptyset$, as for every $\delta>0$ the set $\delta\text{-}\argmin f$ is non-empty, but $\ball_\epsilon(\emptyset)=\emptyset$.
\begin{theorem}\label{thm:alt:convOfMini}
	Let $f,f_n \colon \MS\to\R$. Let $(\epsilon_n)_{n\in\N}\subset [0,\infty)$ with $\epsilon_n \xrightarrow{n\to\infty}0$. 
	Assume $f$ has approachable minimizers, $f_n \xrightarrow{n\to\infty}_{\ubs} f$, and $(\epsilon_n\text{-}\argmin f_n)_{n\in\N}$ is eventually bounded.
	Then
	\begin{equation*}
		d_\subset(\epsilon_n\text{-}\argmin f_n, \argmin f)\xrightarrow{n\to\infty}
		0
	\end{equation*}
	and
	\begin{equation*}
		\inf f_n \xrightarrow{n\to\infty} \inf f
		\eqfs
	\end{equation*}
\end{theorem}
\begin{proof}
		Let $\epsilon >0$. As $f$ has approachable minimizers, there is $\delta > 0$ such that $(3\delta)\text{-}\argmin f \subset \ball_\epsilon(\argmin f)$; also $\argmin f \neq \emptyset$. Let $y\in\argmin f$. As $f_n(y) \xrightarrow{n\to\infty} f(y)$, there is $n_1\in\N$ such that $\inf f_n \leq \inf f + \delta$ for all $n \geq n_1$.		
		As $\epsilon_n \xrightarrow{n\to\infty} 0$, there is $n_2\in\N$ such that $\epsilon_n \leq \delta$ for all $n \geq n_2$.
		As $(\epsilon_n\text{-}\argmin f_n)_{n\in\N}$ is eventually bounded, there is $n_3\in\N$ such that $\diam(B) < \infty$ for $B = \bigcup_{n\geq n_3} \epsilon_n\text{-}\argmin f_n$.
		As $f_n \xrightarrow{n\to\infty}_{\ubs} f$ there is $n_4$ such that $\sup_{x\in B}\abs{f_n(x) - f(x)} \leq \delta$.
		Let $n \geq \max(n_1, n_2, n_3, n_4)$ and $x \in \epsilon_n\text{-}\argmin f_n$. Then
		\begin{equation*}
			f(x) \leq f_n(x) + \delta \leq \inf f_n + 2\delta \leq \inf f + 3\delta\eqfs
		\end{equation*}
		Thus, $x \in (3\delta)\text{-}\argmin f$. By the choice of $\epsilon$ and $\delta$, we obtain $\epsilon_n\text{-}\argmin f_n \subset \ball_\epsilon(\argmin f)$ or equivalently $d_{\subset}(\epsilon_n\text{-}\argmin f_n, \argmin f) \leq \epsilon$.
		
		Finally, we show the convergence of the infima. We already know $\inf f_n \leq \inf f + \epsilon$ for all $\epsilon > 0$ and $n$ large enough.
		If $\inf f_n \xrightarrow{n\to\infty} \inf f$ does not hold, there is a sequence $x_n \in \epsilon_n\text{-}\argmin f_n$ and $\epsilon > 0$ such that $f_n(x_n) < \inf f - \epsilon$ for all $n$ large enough. As before, because of eventual boundedness and uniform convergence on bounded sets, we have $\sup_{k\in\N}\abs{f_n(x_k) - f(x_k)} \xrightarrow{n\to\infty} 0$. Therefore, for all $\epsilon >0$ we have $f(x_n) \leq f_n(x_n) + \epsilon$ for $n$ large enough, which contradicts $f_n(x_n) < \inf f - \epsilon$.
\end{proof}
In the following, we construct examples to show that none of the conditions for one-sided Hausdorff convergence can be dropped.
\begin{example}\label{exa:alt:necessary}
\mbox{ }
\begin{enumerate}[label=(\roman*)]
	\item 
	Let $f,f_n\colon \N_0 \to \R$, $f_n = 1-\indOf{\cb{0,n}}$, $f = 1-\indOf{\cb{0}}$, $d(i,j) = 1$ for $i\neq j$.
	It holds that $f$ is continuous and has approachable minimizers, and the sequence of nonempty sets $\argmin f_n = \cb{0,n}$ is eventually bounded, as $\diam(A)\leq1$ for every $A\subset\N_0$. Furthermore, $f_n$ converges to $f$ uniformly on compact sets, which are exactly the finite subsets of $\N_0$, but not uniformly on bounded sets like $\N_0$ itself.
	There is a subsequence of minimizers $x_n=n\in \argmin f_n$ that is always bounded away from $0$, the minimizer of $f$.
	This shows that uniform convergence on compact sets (instead of bounded sets) is not enough.
	\item
	As above, let $f, f_n\colon \N_0 \to \R$, $f_n = 1-\indOf{\cb{0,n}}$, $f = 1-\indOf{\cb{0}}$, but define $d(i,j) = |i-j|$.
	It holds that $f$ is continuous and has approachable minimizers, and $f_n\ubsconv{n}f$, but the sequence of nonempty sets $\argmin f_n = \cb{0,n}$ is not eventually bounded. 
	Again, there is a subsequence of minimizers $x_n=n\in \argmin f_n$ that is always bounded away from $0$, the minimizer of $f$.
	This shows that eventual boundedness of minimizer sets cannot be dropped.
	\item
	Let $f,f_n\colon \N_0 \to \R$, $f(0) = 0$, $f(i) = \frac1i$, $f_n(i) = f(i)\indOfEvent{i < n}$, and set $d(i,j) = 1$  for $i\neq j$.
	It holds that $f$ is continuous, but $f$ does not have approachable minimizers. The sequence of nonempty sets $\argmin f_n = \cb{0,n, n+1, \dots}$ is eventually bounded and $f_n\ubsconv{n}f$.
	There is a subsequence of minimizers $x_n=n\in \argmin f_n$ that is always bounded away from $0$, the minimizer of $f$.
	This shows that approachability of minimizers of $f$ cannot be dropped.
\end{enumerate}
\end{example}
\subsection{Strong Laws for $\mf c$-Fréchet Mean Sets}\label{ssec:alt:gen}
Let $(\mc Q, d)$ be a metric space, the descriptor space. Let $\mc Y$ be a set, the data space. Let $\mf c\colon \mc Y \times\mc Q \to \R$ be a function, the cost function. Let $(\Omega, \Sigma, \Pr)$ be a probability space that is silently underlying all random variables in this section. Let $Y$ be a random variable with values in $\mc Y$. Denote the $\mf c$-Fréchet mean set of $Y$ as $M = \argmin_{q\in\mc Q} \Ex{\mf c(Y, q)}$.
Let $Y_1, \dots, Y_n$ be independent random variables with the same distribution as $Y$.
Choose $(\epsilon_n)_{n\in\N}\subset [0,\infty)$ with $\epsilon_n \xrightarrow{n\to\infty}0$. Set $M_n = \epsilon_n\text{-}\argmin_{q\in\mc Q} \frac1n \sum_{i=1}^n \mf c(Y_i, q)$.
\begin{assumptions}\mbox{ }
\begin{itemize}
	\item \textsc{Continuity}:  The function $q \mapsto \mf c(Y, q)$ is continuous almost surely.
	\end{itemize}
\end{assumptions}
\begin{theorem}\label{thm:alt:consistency}
	Assume \textsc{HeineBorel}, \textsc{Continuity}, \textsc{UpperBound}, and \textsc{LowerBound}. Then
	\begin{equation*}
		d_\subset(M_n, M) \xrightarrow{n\to\infty}_{\ms{a.s.}} 0
		\eqfs
	\end{equation*}
\end{theorem}
\begin{proof}
Define $F(q) = \Ex{\mf c(Y, q)}$, $F_n(q) = \frac1n\sum_{i=1}^n \mf c(Y_i, q)$.  The proof consists of following steps:
\begin{enumerate}
\item Show that $F_n \xrightarrow{n\to\infty}_{\ubs} F$ almost surely.
\item Reduction to a bounded set.
\item Show that $F$ has approachable minimizers.
\item Show that $M_n$ is eventually bounded.
\item Apply \autoref{thm:alt:convOfMini}.
\end{enumerate}
\underline{\smash{Step 1.}}
To show uniform convergence on bounded sets, we will use the uniform law of large numbers, \autoref{thm:ULLN} below.
Let $B\subset\mc Q$ be a bounded set.
By \textsc{HeineBorel}, $\overline{B}$ is compact.
By \textsc{Continuity}, $q \mapsto \mf c(Y, q)$ is almost surely continuous.
By \textsc{UpperBound}, $\Ex{ \sup_{q\in B}\abs{\mf c(Y, q)}} < \infty$. Thus, \autoref{thm:ULLN} implies that $q \mapsto F(q)$ is continuous and
\begin{equation*}
	\sup_{q\in B}  \abs{F_n(q)-F(q)} \xrightarrow{n\to\infty}_{\ms{a.s.}}0
	\eqfs
\end{equation*}
Fix an arbitrary element $o\in\mc Q$. For all bounded sets $B$, there is $\delta \in\N$ such that $B \subset \ball_\delta(o)$. By the previous considerations, uniform convergence holds almost surely for all $(\ball_{\delta}(o))_{\delta\in\N}$.
Thus, $F_n \xrightarrow{n\to\infty}_{\ubs} F$ almost surely.

\noindent
\underline{\smash{Step 2.}}
Find $B_1\subset\mc Q$ and a random variable $N_1\in\N$ as in step 2 in the proof of \autoref{thm:consistency}.

\noindent
\underline{\smash{Step 3.}}
Clearly, $M \subset B_1$ is bounded. Furthermore, for all $\epsilon > 0$ small enough the set $D_\epsilon = \overline{B_1 \setminus \ball_\epsilon(M)}$ is not empty (if it is, increase $\delta$), does not contain any element of $M$ and, by \textsc{HeineBorel}, is compact. Thus, the continuous function $q\mapsto F(q)$ attains its infimum on $D_\epsilon$ where $\inf_{q\in D_\epsilon} F(q) > \inf_{q\in \mc Q} F(q)$. Take $\zeta = \min(1, \frac12(\inf_{q\in D_\epsilon} F(q) - \inf_{q\in \mc Q} F(q)))$. Then $\zeta\text{-}\argmin_{q\in\mc Q} F(q) \subset \ball_\epsilon(M)$, i.e., $F$ has approachable minimizers. 

\noindent
\underline{\smash{Step 4.}}
For $\epsilon_n < 1$ and $n \geq N_1$, it holds $M_n \subset B_1$. Thus, $(M_n)_{n\in\N}$ is eventually bounded almost surely.

\noindent
\underline{\smash{Step 5.}}
Finally, \autoref{thm:alt:convOfMini} implies $d_\subset(M_n, M) \xrightarrow{n\to\infty}_{\ms{a.s.}} 0$.
\end{proof}
 \section{Auxiliary Results}\label{sec:aux}
There are many versions of uniform laws of large numbers in the literature. We state and prove one version that is tailored to our needs.
\begin{theorem}\label{thm:ULLN}
Let $(\mc Y, \Sigma_{\mc Y})$ be a measurable space and $Y$ be a random variable with values in $\mc Y$. Let $Y_1, \dots, Y_n$ be independent and have the same distribution as $Y$.
Let $(\mc Q, d)$ be a metric space and $B\subset \mc Q$ compact. 
Let $f \colon \mc Y \times B \to \R$ be such that $q \mapsto f(Y, q)$  is almost surely continuous. Assume there is a random variable $Z$ such that $\abs{f(Y, q)} \leq Z$ for all $q \in B$ with $\Ex{Z} < \infty$. Then $q \mapsto \Ex{f(Y, q)}$ is continuous and 
\begin{equation*}
	\sup_{q\in B} \abs{\frac 1n \sum_{i=1}^n f(Y_i, q) - \Ex{f(Y, q)}} \xrightarrow{n\to\infty}_{\ms{a.s.}} 0 
	\eqfs
\end{equation*}
\end{theorem}
\begin{proof}
Let $\epsilon > 0$. As $B$ is compact, there is a finite set $\cb{q_1, \dots, q_k} \subset \mc Q$ such that $B \subset \bigcup_{\ell=1}^k\ball_\epsilon(q_\ell)$. We split the supremum,
\begin{align*}
	&\sup_{q\in B} \abs{\frac 1n \sum_{i=1}^n f(Y_i, q) - \Ex{f(Y, q)}}
	\\&\leq
	\sup_{\ell \in \cb{1,\dots, k}} 
	\sup_{q\in \ball_\epsilon(q_\ell)}  
		\abs{\frac 1n \sum_{i=1}^n \br{f(Y_i, q) - f(Y_i, q_\ell)} - \Ex{f(Y, q) - f(Y, q_\ell)}} 
	\\&\quad+ 
	\sup_{\ell \in \cb{1,\dots, k}} \abs{\frac 1n \sum_{i=1}^n f(Y_i, q_\ell) - \Ex{f(Y, q_\ell)}}
	\eqfs
\end{align*}
For the second summand, by the standard strong law of large numbers applied to each $\ell\in\{1,\dots,k\}$ with $\Ex{Z} < \infty$, 
\begin{align*}
\sup_{\ell \in \cb{1,\dots, k}} \abs{\frac 1n \sum_{i=1}^n f(Y_i, q_\ell) - \Ex{f(Y, q_\ell)}} \xrightarrow{n \to \infty}_{\ms{a.s.}} 0
\eqfs
\end{align*}
For the first summand,
\begin{align*}
	&\sup_{\ell \in \cb{1,\dots, k}} 
	\sup_{q\in \ball_\epsilon(q_\ell)}  
		\abs{\frac 1n \sum_{i=1}^n \br{f(Y_i, q) - f(Y_i, q_\ell)} - \Ex{f(Y, q) - f(Y, q_\ell)}} 
	\\& \leq
	\frac 1n \sum_{i=1}^n \sup_{q,p\in B,\, \sol qp \leq \epsilon} \abs{f(Y_i, q) - f(Y_i, p)} 
	+ 
	\Ex*{\sup_{q,p\in B,\, \sol qp \leq \epsilon} \abs{f(Y, q) - f(Y, p)}} 
	\eqfs
\end{align*}
By the standard strong law of large numbers with $\Ex{Z} < \infty$, 
\begin{align*}
	\frac 1n \sum_{i=1}^n \sup_{q,p\in B,\, \sol qp \leq \epsilon} \abs{f(Y_i, q) - f(Y_i, p)} 
	&\xrightarrow{n \to \infty}_{\ms{a.s.}} 
	\Ex*{\sup_{q,p\in B,\, \sol qp \leq \epsilon} \abs{f(Y, q) - f(Y, p)} }
	\eqfs
\end{align*}
Thus,
\begin{equation}\label{eq:ulln:as_limsup}
	\PrOf{
		\limsup_{n\to\infty} \sup_{q\in B} \abs{\frac 1n \sum_{i=1}^n f(Y_i, q) - \Ex{f(Y, q)}}
		\leq
		a_\epsilon
	} = 1
	\eqcm
\end{equation}
where $a_\epsilon = 2 \Ex*{\sup_{q,p\in B,\, \sol qp \leq \epsilon} \abs{f(Y, q) - f(Y, p)} }$.
As $q \mapsto f(Y, q)$ is almost surely continuous and $B$ is compact, $q \mapsto f(Y, q)$ is almost surely uniformly continuous, i.e., for all $\delta > 0$ there is $\varepsilon(\delta, Y) > 0$ such that $\abs{f(Y,q)-f(Y,p)} \leq \delta$ for all $\ol qp \leq \varepsilon(\delta, Y)$.
As $\Ex{Z} < \infty$, we can use dominated convergence to obtain 
\begin{align*}
	\lim_{\epsilon\searrow0} 	\Ex*{\sup_{q,p\in B,\, \sol qp \leq \epsilon} \abs{f(Y, p) - f(Y, p)} }
	&=
	\Ex*{\lim_{\epsilon\searrow0}\sup_{q,p\in B,\, \sol qp \leq \epsilon} \abs{f(Y, p) - f(Y, p)} }
	=
	0
	\eqfs
\end{align*}
Thus, $a_\epsilon \xrightarrow{\epsilon\searrow0} 0$. Together with \eqref{eq:ulln:as_limsup}, this implies
\begin{equation*}
 \sup_{q\in B} \abs{\frac 1n \sum_{i=1}^n f(Y_i, q) - \Ex{f(Y, q)}} \xrightarrow{n \to \infty}_{\ms{a.s.}} 0
 \eqfs
\end{equation*}
We have also shown that $q\mapsto \Ex{f(Y, q)}$ is continuous, as $\abs{\Ex{f(Y, q)}-\Ex{f(Y, p)}} \leq a_{\sol qp}$.
\end{proof}
 \begin{lemma}\label{lmm:nondec}
Let $h \colon [0,\infty) \to [0,\infty)$ be a non-decreasing function.
Define $H \colon [0,\infty) \to [0,\infty), x\mapsto \int_0^x h(t) \dl t$.
Let $x, y \geq 0$. Then
\begin{enumerate}[label=(\roman*)]
\item $\abs{H(x) - H(y)} \leq \abs{x-y}h(\max(x, y))$.
\end{enumerate}
Assume, there is $b \in [1,\infty)$ such that $h(2u) \leq b h(u)$ for all $u \geq 0$. Then
\begin{enumerate}[label=(\roman*),start=2]
\item $\frac12 h(x) + \frac12 h(y) \leq h (x+y) \leq b \br{h(x) + h(y)}$,
\item $H(\abs{x-y}) - H(x) \geq b^{-1} H(y) - 2 y h(x)$.
\end{enumerate}
\end{lemma}
\begin{proof}
\begin{enumerate}[label=(\roman*)]
\item This is a direct consequence of the mean value theorem.
\item As $h$ is non-decreasing, $\max(h(x), h(y)) \leq h(x+y) \leq \max(h(2x), h(2y))$.
	By the definition of $b$ and with $\frac12(u+v) \leq \max(u,v) \leq u+v$ for $u,v \geq 0$ the claim follows.
\item 
	First, consider the case $x\geq y$.
 	Define $f(x,y) = H(x-y)-H(x)- b^{-1} H(y) + 2 y h(x)$.
 	We want to show $f(x, y)\geq0$.
	The derivative of $f$ with respect to $y$ is 
	\begin{equation*}
		\partial_y f(x, y) = -h(x-y) - b^{-1}h(y) + 2 h(x)\eqfs
	\end{equation*} 
	By applying the first inequality of (ii) to $h(x) = h((x-y) + y)$, we obtain $\partial_y f(x, y)  \geq 0$ as $b^{-1} \leq 1$.
	Hence, $f(x, y) \geq f(x, 0) = 0$, as $H(y) = 0$.
	
	Now, consider the case $x \leq y$. Set $g(x,y) = H(y-x) - H(x) - b^{-1} H(y) + 2 y h(x)$, which yields
	\begin{equation*}
		\partial_y g(x, y) = h(y-x) - b^{-1} h(y) + 2 h(x)\eqfs
	\end{equation*} 
	By applying the second inequality of (ii) to $h(y) = h((y-x) + x)$, we obtain $\partial_y g(x, y)  \geq 0$ as $b^{-1} \leq 1$.
	Thus, $g(x, y) \geq g(x, x) = -(1 + b^{-1}) H(x) + 2 x h(x)$ as $H(0) = 0$.
	By the definition of $H$, as $h$ is non-decreasing, $H(x) \leq x h(x)$. Hence, $g(x, y) \geq 0$ as $1 + b^{-1} \leq 2$.

	Together, we have shown $H(\abs{x-y}) - H(x) - b^{-1} H(y) + 2 y h(x) \geq 0$ for all $x,y\geq0$.
\end{enumerate}
\end{proof}
\begin{lemma}\label{lmm:nondec:existence}
	Let $X$ be a random variable with values in $[0,\infty)$. Then there is a strictly increasing, continuous, and concave function $h \colon [0,\infty) \to [0,\infty)$ with $h(\delta)\xrightarrow{\delta\to\infty}\infty$ such that $\Ex{h(X)} < \infty$.
\end{lemma}
\begin{proof}
	If there is $K > 0$ such that $\Prof{X < K} = 1$ take $h(x) = x$. Now, assume that $X$ is not almost surely bounded.
	We first construct a non-decreasing function $\tilde h \colon [0,\infty) \to [0,\infty)$ such that $\tilde h (x) \xrightarrow{x\to\infty}\infty$ with $\Ex{\tilde h(X)} < \infty$.
	Then we construct a function $h$ from $\tilde h$ with all desired properties.
	
	Let $F$ be the distribution function of $X$, $F(x) = \Prof{X \leq x}$.
	Let $z_1 = 0$ and $z_{n+1} = \inf \cbOf{x\geq z_n +1 \,\big\vert\, 1-F(x) \leq \frac1n}$. As $F(x)\xrightarrow{x\to\infty} 1$, $z_n < \infty$. Furthermore, $z_{n+1}-z_n \geq 1$. Moreover, as $X$ is not almost surely bounded, $1-F(x) > 0$ for all $x\geq 0$. Set 
	\begin{align*}
		g(x) &= \sum_{n=1}^\infty (z_{n+1}-z_n)^{-1}n^{-2}\ind_{[z_n, z_{n+1})} (x)
		\eqcm
		\\
		\tilde h(x) &= \int_0^x \frac{g(t)}{1-F(t)} \dl t
		\eqfs
	\end{align*}
	Then
	\begin{align*}
		\lim_{x\to\infty} \tilde h(x) 
		= 
		\int_0^\infty \frac{g(t)}{1-F(t)} \dl t
		\geq
		\sum_{n=1}^\infty n^{-1} 
		= 
		\infty
		\eqfs
	\end{align*}
	Moreover, $\tilde h(x)$ is strictly increasing, as $g(t)\geq 0$ and $1-F(t) \geq 0$.
	The function $\tilde h$ is continuously differentiable everywhere except at point $z_n$, $n\in\N$. Thus,
	\begin{align*}
		\Ex{\tilde h(X)}
		&=
		\int_0^\infty \PrOf{\tilde h(X) > t} \dl t
		\\&=
		\int_0^\infty \PrOf{X > \tilde h^{-1}(t)} \dl t
		\\&=
		\int_0^\infty \tilde h\pr(t) \PrOf{X > t} \dl t
		\\&=
		\int_0^\infty g(t) \dl t
		\\&=
		\sum_{n=1}^\infty n^{-2}
		<
		\infty
		\eqfs
	\end{align*}
	Let $a_0 = 1$, $x_0=0$, $x_{n+1} = \inf\cbOf{x \geq x_n + a_n^{-1} \,\big\vert\, \tilde h(x) \geq n+1}$ and $a_{n+1} = (x_{n+1} - x_n)^{-1}$. Let $h \colon [0,\infty) \to [0,\infty)$ be the linear interpolation of $(x_n, n)_{n\in\N_0}$. 
	As $\tilde h(x) \xrightarrow{x\to\infty}\infty$, all $x_n$ are finite. Hence, $h(x) \xrightarrow{x\to\infty}\infty$. 
	Because of $a_n > 0$, $h$ is strictly increasing.
	Furthermore, $a_{n+1} \leq a_n$ as $x_{n+1} \geq x_n + a_n^{-1}$. As $h$ is continuous and $a_n$ is the derivative of $h$ in the interval $(x_n, x_{n+1})$, $h$ is concave. 
	Lastly, $h(x) \leq \tilde h(x) + 1$. Thus, $\Ex{h(X)} < \infty$.
\end{proof}
 \begin{lemma}[Fréchet means in product spaces]\label{lmm:products}
Let $K\in\N$. Let $(\mc Q_1, d_1), \dots, (\mc Q_K, d_K)$ be metric spaces. 
Let $\alpha\geq1$.
Set $\mc Q := \bigtimes_{k=1}^K \mc  Q_k$ and $d \colon \mc Q \times \mc Q \to [0,\infty)$, $d(q,p) := (\sum_{k=1}^K d_k(q_k, p_k)^\alpha)^\frac1\alpha$. Then $(\mc Q, d)$ is a metric space. Let $Y = (Y^1, \dots, Y^K)$ be a tuple of random variables such that $Y^k$ has values in $\mc Q_k$ and $\Ex{d(Y,o)^{\alpha-1}} < \infty$ for an element $o\in\mc Q$. Let $M$ be the $\alpha$-Fréchet mean set of $Y$ in $(\mc Q, d)$, and $M^k$ be the $\alpha$-Fréchet mean set of $Y^k$ in $(\mc Q_k, d_k)$. Then $M$ is the Cartesian product of the sets $M^k$, i.e.,
\begin{equation*}
	M = \bigtimes_{k=1}^K M^k
	\eqfs
\end{equation*}
\end{lemma}
\autoref{lmm:products} can be proven by straight forward calculations.
 \begin{lemma}\label{lmm:precompact}
Let $(\mc Q, d)$ be a complete metric space and $A\subset \mc Q$. Assume that all closed subsets $B\subset A$ are compact. Then $\overline A$ is compact.
\end{lemma}
\begin{proof}
	Let $(a_k)_{k\in\N}\subset \overline{A}$. We show that $(a_k)_{k\in\N}$ has a converging subsequence with limit $a \in \overline{A}$, which implies compactness of $\overline{A}$.
	
	Let $\delta_n = 2^{-n}$. Define $B_n := \mc Q \setminus \ball_{\delta_n}(\mc Q \setminus A)$. The sets $B_n$ are closed and subsets of $A$. Thus, they are compact. Let $b_k^n\in\argmin_{b \in B_n} d(b, a_k)$. Such an element exists as $B_n$ is compact. Furthermore, $d(b_k^n, a_k) \leq \delta_n$. Define the subindex sequences $(k(n,\ell))_{\ell\in\N} \subset \N$ such that $k(0, \ell) = \ell$ and $(b^n_{k(n, \ell)})_{\ell\in\N}$ is a converging subsequence of $(b^n_{k(n-1, \ell)})_{\ell\in\N}$ with limit $b^n_{k(n, \ell)} \xrightarrow{\ell\to\infty}b^n_\infty$ and $d(b^n_{k(n, \ell)}, b^n_\infty) \leq \delta_\ell$. By the triangle inequality $d(b^{n_1}_k, b^{n_2}_k) \leq d(b^{n_1}_k, a_k) + d(b^{n_2}_k, a_k) \leq \delta_{n_1}+\delta_{n_2}$. Thus, $d(b^{n_1}_\infty, b^{n_2}_\infty) \leq \delta_{n_1}+\delta_{n_2}$, which makes $(b^{n}_\infty)_{n\in\N}$ a Cauchy-sequence. Define $a$ as its limit, i.e., $b^{n}_\infty \xrightarrow{n\to\infty} a$. As $b^{n}_\infty\in B_n \subset\overline{A}$, also $a\in\overline{A}$. Finally, the triangle inequality yields
	\begin{equation*}
		d(a_{k(n, n)}, a) \leq d(a_{k(n, n)}, b^n_{k(n, n)}) + d(b^n_{k(n, n)}, b^{n}_\infty) + d(b^{n}_\infty, a) \xrightarrow{n\to\infty}0\eqcm
	\end{equation*}
	i.e., $(a_{k(n, n)})_{n\in\N}$ is subsequence of $(a_k)_{k\in\N}$ which converges in $\overline{A}$.
\end{proof}
 \end{appendix}
\bibliographystyle{apalike}
\bibliography{literature}
\end{document}